% ---------------------------------------------------------
\chapter{Zusammenfassung und Ausblick}\label{chap:summary}
% ---------------------------------------------------------

Hier kommt die Gaußsche Summenformel
\[
\sum_{i=1}^n i = \frac{n(n+1)}{2}.
\]


\begin{figure}[th]
  \centering
\tikz{\pie [explode = {0, 0, 0.2, 0} ] {
  21/ Partei A,
  11/ Partei B,
  42/ Partei C,
  26/ Partei D }}
\caption{Ergbnisse der Foobar-Wahl 2020.}\label{fig:wahl2020}
\end{figure}

In Abbildung \ref{fig:wahl2020} sehen Sie die Wahlergebnisse.


Mehr zu \LaTeX\xspace erfahren Sie unter
\url{https://en.wikibooks.org/wiki/LaTeX}.
